\section{原创性与二次开发说明}

\subsection{业务设计与代码来源}

一纸成课为本赛事独立设计和实现的作品，不是对现有错题本、扫描、出题或在线课程项目的 fork，也没有在某个开源应用上更换名称和界面后提交。纸面到卡片的业务路径、七 Agent 分工、命题与质检反馈、三级验算语义、卡片修订历史、事件回放、函数图交互和知识谱系统计均围绕本作品需求开发。通用 Web 框架与符号计算库提供基础设施，不等同于复用已有同类作品。

开发过程中使用 AI 编码助手讨论方案、生成部分代码草稿、补充测试用例和开展独立评审。开发者负责确定 Must/Cut 范围、审查关键逻辑、选择是否采纳建议，并执行实际验收。架构上否决了超出当前需求的状态机框架、向量库和绘图库；安全上没有采用黑名单过滤建议，而是重新实现白名单解析；产品上没有把 mock 结果描述为真实识别。关键采纳与否决记录保存在 \file{docs/ai-logs/}，可以与提交历史及测试变化对照。

项目没有训练或微调自有基础模型。智谱、DeepSeek 以及配置中支持的其他 Provider 属于第三方模型与服务，运行时调用的模型能力不列为自研成果。自写 mock 样例用于复现和测试，三套教学内容不包含真实学生资料。BGE-M3 只在可选混合检索评测中作为现成嵌入模型使用，也不属于本作品训练的模型。

\subsection{第三方组件的边界}

FastAPI、React、sympy、PyMuPDF、PyInstaller、Electron、KaTeX、python-docx 等组件分别提供 Web 服务、界面、符号计算、文档处理、打包、桌面壳和公式渲染能力。它们按照各自许可证使用，依赖名称、用途和许可在附录 B 列出。仓库根目录的 \file{THIRD_PARTY_NOTICES.md} 与桌面资源中的许可证文本用于说明再分发义务。

源码采用 MIT 许可，但桌面二进制包含 PyMuPDF 等带有 AGPL 义务的组件。根许可证不能被描述为把所有组件都变成 MIT。公开源码、对应获取方式和 NOTICE 是当前交付采取的措施，若未来分发渠道要求整包采用另一种许可结构，需要更换组件或取得适当商业授权。该问题属于许可合规边界，不影响对业务代码原创来源的说明，但也不能被省略。

\begin{table}[H]
\centering\footnotesize
\caption{原创性核验对象与判断方式}
\begin{tabularx}{\textwidth}{@{}L{3.0cm}Y L{3.7cm}@{}}
\toprule
\textbf{对象} & \textbf{来源说明} & \textbf{核验方式} \\
\midrule
业务流程与数据结构 & 围绕纸面资料转学习动作独立设计，无同类项目 fork & 对照仓库初始设计、数据表和后续提交 \\
七 Agent 与反馈回路 & 按能力异构和独立否决需求实现，编排器为本项目代码 & 查看 \file{orchestrator.py}、Agent 文件和流水线测试 \\
前端交互与图形 & 单页六面板、事件看板、SVG 函数图和谱系为本项目实现 & 查看组件源码、前端纯函数测试与运行截图 \\
AI 辅助内容 & 草稿、候选用例和评审建议由 AI 参与，人工判断后落地 & 查看快照中的提示词、否决理由与对应提交 \\
教学样例与评测集 & 二次函数、一元二次方程、三角函数样例自写，评测为构造集 & 查看 fixtures、样例生成脚本和结果 JSON \\
基础库与模型 & 使用通用开源组件及第三方 API，不声称这些组件和模型自研 & 查看依赖锁文件、NOTICE、Provider 配置与附录 B \\
早期草稿与内部评审 & 未作为公开成果发表，不授权其他参赛作品重复使用 & 公开仓库只含本次交付所需材料，历史不补造 \\
\bottomrule
\end{tabularx}
\end{table}

原创性不等于每一行基础代码都由手工独立完成，也不等于没有使用通用库和 AI。这里的声明重点是：没有套用他人同类作品，业务设计与工程取舍可解释，第三方和 AI 的作用可核对，关键实现能够运行和测试。公开仓库的细粒度历史用于展示过程，不能单独作为原创性的唯一证明；源码、快照、测试和实际运行共同构成核验材料。
